\documentclass[11pt,a4paper]{article}
\usepackage[T1]{fontenc}
\usepackage[utf8]{inputenc}
\usepackage{lmodern}
\usepackage[margin=1in]{geometry}
\usepackage{amsmath,amssymb,amsthm,mathtools}
\usepackage{microtype}
\usepackage{booktabs,array,tabularx}
\usepackage[numbers,sort&compress]{natbib}
\usepackage{xcolor}
\usepackage{hyperref}
\definecolor{linkblue}{HTML}{185C78}
\hypersetup{colorlinks=true,linkcolor=linkblue,citecolor=linkblue,
  urlcolor=linkblue,pdftitle={Optimal Communication-Regret Tradeoffs for Federated Two-Point Bandit Optimization},
  pdfauthor={Colombo Team}}
\setlength{\parindent}{1.2em}
\setlength{\parskip}{3pt}
\setlength{\emergencystretch}{2em}
\linespread{1.03}
\numberwithin{equation}{section}
\newtheorem{theorem}{Theorem}[section]
\newtheorem{lemma}[theorem]{Lemma}
\newtheorem{proposition}[theorem]{Proposition}
\newtheorem{corollary}[theorem]{Corollary}
\theoremstyle{remark}
\newtheorem{remark}[theorem]{Remark}
\newcommand{\R}{\mathbb R}
\newcommand{\E}{\mathbb E}
\newcommand{\calR}{\mathfrak R}
\newcommand{\norm}[1]{\left\lVert #1\right\rVert}
\newcommand{\ip}[2]{\left\langle #1,#2\right\rangle}
\newcommand{\B}{\mathbb B_2}
\newcommand{\Sph}{\mathbb S^{d-1}}
\newcommand{\ind}{\mathbf 1}
\DeclareMathOperator{\spanop}{span}
\DeclareMathOperator{\Unif}{Unif}
\DeclareMathOperator{\argmin}{argmin}
\newcolumntype{Y}{>{\raggedright\arraybackslash}X}
\title{Optimal Communication--Regret Tradeoffs for\\
Federated Two-Point Bandit Optimization}
\author{Colombo Team\\[3pt]
  \small\href{mailto:colombo.2026@outlook.com}{\texttt{colombo.2026@outlook.com}}}
\date{}

\begin{document}
\maketitle

\begin{abstract}
We determine the minimax regret of federated two-point bandit convex
optimization with intermittent communication. There are $M$ workers,
$T=KR$ time steps, and a communication round after every $K$ steps.
At each step, each worker selects two query points simultaneously and
incurs the average of their losses. The losses are differentiable,
convex, and globally $G$-Lipschitz, and may be chosen by a nonanticipating
adaptive adversary. Actions are unrestricted, while the workers share a
Euclidean-ball comparator of radius $B$. We prove that the minimax
average expected regret is
\[
\Theta\!\left(GB\min\left\{1,\sqrt{\frac dT},
                 \sqrt{\frac1R+\frac d{MT}}\right\}\right),
\]
uniformly over all positive integer parameters. The upper bound combines
independent two-point learning, communication-block batching, and the
zero action. Its analysis explicitly controls the random hindsight
comparator under adaptive losses. The lower bound separates two sources
of difficulty: information that remains local between communications,
and information that remains unobserved even after pooling all workers'
feedback. Both constructions use residuals of fresh spherical directions;
a fixed differentiable radial penalty makes them valid for arbitrary
unbounded action distributions. The resulting characterization identifies
when collaboration helps and gives matching bounds throughout the
communication, dimension, and constant-regret regimes.
\end{abstract}

\noindent\textbf{Keywords.}
Bandit convex optimization; federated online learning; intermittent
communication; minimax regret; adaptive adversaries.

\begin{samepage}
\noindent\textbf{AI Usage.}
Nearly the entire research pipeline for this paper was carried out by
Colombo, powered by GPT-5.6 Sol.
\par
\end{samepage}

\section{Introduction}
\label{sec:intro}

When several learners observe only function values, communication can
combine complementary observations held by different workers. This benefit
must be weighed against the time during which local observations cannot
be shared. In an online problem, the tradeoff is especially delicate:
losses continue to arrive between communications, and an adaptive
adversary can respond to the learners' earlier decisions. The relevant
question is therefore not only how many evaluations are available, but
also when their information becomes available to each worker.

We study this question in federated bandit convex optimization. Each of
$M$ workers makes two simultaneous evaluations per time step, and all
workers may exchange their histories once every $K$ steps. The horizon
is $T=KR$, where $R$ is the number of communication rounds. Both
evaluations incur loss, and performance is measured against one common
point chosen in hindsight from a Euclidean ball. Losses may differ
arbitrarily across workers and may depend on the entire past interaction.
The comparator is bounded, but the query points are not.

\paragraph{Main result.}
For globally $G$-Lipschitz, differentiable convex losses in $\R^d$ and
comparator radius $B$, the optimal average expected regret is
\[
GB\sqrt{\frac{\min\{T,d,K+d/M\}}{T}}
\]
up to universal constants. This formula holds for every positive integer
$M,K,R,d$. It combines the independent-learning scale $GB\sqrt{d/T}$,
the collaborative scale $GB\sqrt{1/R+d/(MT)}$, and the trivial cap $GB$.
In particular, independent learning is optimal when $d\le K$, the number
of communication rounds controls the rate when $K<d<MK$, and pooled
bandit information determines the rate when $d\ge MK$, until the regret
reaches constant scale.

\paragraph{Why the lower bound depends on communication.}
In our lower-bound family, each worker's two values reveal at most two
adaptive linear measurements per step.
Our communication adversary draws a common hidden direction and, at
each worker, removes only the component already determined by that
worker's own queries. Every charged linear loss then has conditional
mean zero, while the residuals across workers retain a substantial
component in the common hidden direction. Restarting this construction
before a communication or after order-$d$ local observations produces
the communication lower bound. A second adversary removes the subspace
measured by \emph{all} workers, yielding the dimension lower bound even
with communication after every step. Taking the larger of the two
bounds recovers the complete minimax rate.

\paragraph{Analytical issues.}
Two features of the model require care. First, under an adaptive
adversary, the best comparator in hindsight is random and correlated
with the gradient-estimation noise. The upper bound keeps this
comparator inside a pathwise inequality and controls the resulting
terminal martingale term. Second, unrestricted actions can have
nonintegrable norms. A common radial penalty, zero on the comparator
ball, makes every lower-bound loss uniformly bounded below without
restricting the learner's actions. These arguments allow arbitrary
measurable randomized protocols, including shared randomization.

\subsection{Prior work and the comparison in this paper}

Multi-point bandit feedback was studied by Agarwal, Dekel, and
Xiao~\cite{ADX10}. Shamir~\cite{Shamir17} analyzed a symmetric two-point
estimator whose second moment is linear in dimension for general
Lipschitz convex losses. We use this standard estimator and projected
online gradient descent~\cite{Zinkevich03}. Their combination supplies
the independent-learning baseline; our proof also treats the random
hindsight comparator explicitly.

Patel, Wang, Saha, and Srebro~\cite{PWSS23} formulated federated online
and bandit convex optimization against adaptive adversaries. Their
Theorem~5.1 gives the following general-Lipschitz guarantee for
\textsc{FedOSGD}, in the same average-regret normalization:
\begin{equation}
 O\!\left(GB\left[\frac1{\sqrt T}+\sqrt{\frac d{MT}}
       +\ind_{\{K>1\}}\frac{d^{1/4}}{\sqrt R}\right]\right).
 \label{eq:prior}
\end{equation}
Their Algorithm~3 uses simultaneous symmetric queries, both of which
incur loss; the comparator is shared and lies in a Euclidean ball.
Their conclusion asks whether the communication terms are tight.
Table~\ref{tab:comparison} aligns this published guarantee with the
bounds proved here. Our theorem determines the optimal communication
dependence for the explicitly simultaneous-query model with unrestricted
heterogeneity. Smoothness-dependent and small-heterogeneity refinements
are outside its scope.

\begin{table}[t]
\centering
\small
\renewcommand{\arraystretch}{1.35}
\setlength{\tabcolsep}{4pt}
\begin{tabularx}{\textwidth}{@{}>{\raggedright\arraybackslash}p{0.23\textwidth}p{0.40\textwidth}Y@{}}
\toprule
Result or strategy & Average regret divided by $GB$ & Conditions and role\\
\midrule
Independent learning
& $O\!\left(\min\{1,\sqrt{d/T}\}\right)$
& No communication; standard two-point estimator.\\
Patel et al.~\cite{PWSS23}\newline Theorem~5.1
& $O\!\left(\dfrac1{\sqrt T}+\sqrt{\dfrac d{MT}}
       +\ind_{\{K>1\}}\dfrac{d^{1/4}}{\sqrt R}\right)$
& General Lipschitz losses; guarantee for local-update \textsc{FedOSGD}.\\
Block batching\newline (this paper)
& $O\!\left(\min\{1,\sqrt{1/R+d/(MT)}\}\right)$
& One common center per communication block.\\
Minimax rate\newline (this paper)
& $\Theta\!\left(\min\{1,\sqrt{d/T},
                 \sqrt{1/R+d/(MT)}\}\right)$
& Matching bounds for all $M,K,R,d\ge1$ in Section~\ref{sec:model}.\\
\bottomrule
\end{tabularx}
\caption{Comparison for $T=KR$, with no smoothness or small-heterogeneity
assumption. Both query losses are charged, and all rates are average
regret. The first and third rows include the zero-action cap. The
published bound in the second row is an algorithmic upper guarantee;
the last row is a minimax statement for the precise information and
timing rules specified below.}
\label{tab:comparison}
\end{table}

For a fair quantitative comparison, combine \eqref{eq:prior} with
independent learning and the zero action. Define the resulting rate
expression
\begin{equation}
 U_{\rm prior}=GB\min\left\{1,\sqrt{\frac dT},
 \frac1{\sqrt T}+\sqrt{\frac d{MT}}
       +\ind_{\{K>1\}}\frac{d^{1/4}}{\sqrt R}\right\}.
 \label{eq:prior-envelope}
\end{equation}
For $M=K=s$, $d=R=s^2$, and $T=s^3$, this expression has order
$GBs^{-1/2}$, whereas our minimax rate is $GBs^{-1}$. This is an
improvement over the compared upper-bound guarantees; it does not assert
a lower bound on the actual performance of \textsc{FedOSGD}.

Batching itself is a standard distributed-learning device. Dekel et
al.~\cite{DGSX12} analyze distributed mini-batches for smooth losses
with stochastic inputs. Woodworth et al.~\cite{WBSS21} characterize
distributed stochastic convex optimization with intermittent
communication and first-order feedback. These works provide useful
context, but their objectives and information models differ from
adaptive online regret with exact function values. The lower bounds
in this paper are consequently proved within the stated bandit model.

\paragraph{Organization.}
Section~\ref{sec:model} defines the model and states the main theorem
and its consequences. Section~\ref{sec:upper} proves the upper bound.
Section~\ref{sec:tools} develops the probability and information lemmas
used in the lower bounds. Sections~\ref{sec:local} and
\ref{sec:global} give the communication and dimension constructions;
Section~\ref{sec:combine} combines them. Section~\ref{sec:discussion}
discusses the scope of the characterization.

\section{Model, main theorem, and consequences}
\label{sec:model}

\subsection{Interaction and information}

Let $M,K,R,d$ be positive integers, let $T=KR$, and initially assume
$G,B>0$. Write $[n]=\{1,\ldots,n\}$, let $\norm{\cdot}$ denote the
Euclidean norm, and set $\B(B)=\{x\in\R^d:\norm{x}\le B\}$. The
interaction consists of $T$ synchronous steps. At step $t$:
\begin{enumerate}
\item A nonanticipating adversary selects convex, differentiable losses
  $f_t^m:\R^d\to\R$, $m\in[M]$, each globally $G$-Lipschitz. Its
  joint choice may depend on all previously played points and returned
  values, together with its own randomness, but not on the current
  queries or unrevealed current random choices of the workers.
\item Worker $m$ selects a pair $(x_t^{m,1},x_t^{m,2})\in(\R^d)^2$.
  All workers' current pairs are selected before any current function
  value is returned. The two points in a pair may be equal.
\item Worker $m$ receives only $f_t^m(x_t^{m,1})$ and
  $f_t^m(x_t^{m,2})$, and both values count toward its loss. At steps
  $K,2K,\ldots,RK$, the workers may then exchange their complete
  histories.
\end{enumerate}
Between communications, a worker has access only to its own history,
its randomization, and information received at earlier communications.
Other workers' actions and observations are not public within a block.
The functions themselves are not revealed. In particular, even when
$K=1$, communication follows the current evaluations.

Private and shared algorithmic randomness are allowed. There are no
restrictions on computation, message length, the span of the queries,
or the query norms. All maps defining protocols and adversarial strategies
are measurable. The
communication schedule is fixed in advance. The final communication
does not affect regret, but is included in the convention $T=KR$.
The parameters are available to the learner.

\subsection{Expected hindsight regret}

For a protocol $A$ and an adversarial strategy $P$, define the sample-path
average regret by
\begin{equation}
 Z_T(A,P)=\frac1{2MT}\left[
 \sum_{t=1}^T\sum_{m=1}^M\sum_{j=1}^2 f_t^m(x_t^{m,j})
 -2\min_{u\in\B(B)}\sum_{t=1}^T\sum_{m=1}^M f_t^m(u)
 \right].
 \label{eq:path-regret}
\end{equation}
Let $\operatorname{Reg}(A,P)=\E Z_T(A,P)$ whenever this extended
expectation exists. If both the positive and negative parts have
infinite expectation, set $\operatorname{Reg}(A,P)=+\infty$;
such a protocol has no finite worst-case guarantee. No moment restriction is imposed
on its actions. This convention makes the minimax value
\begin{equation}
 \calR_{M,K,R,d}(G,B)=\inf_A\sup_P\operatorname{Reg}(A,P)
 \label{eq:model}
\end{equation}
well defined over all measurable protocols.

The expectation is taken \emph{after} forming the complete difference
in \eqref{eq:path-regret}; arbitrary additive constants in the losses
therefore cancel pathwise. The minimum exists by continuity and
compactness and remains inside the expectation. We compare with the
best fixed point for the realized loss sequence, rather than the losses
that an adaptive adversary would have generated along a different
trajectory. Thus \eqref{eq:model} measures expected hindsight regret,
not policy regret. The ball restricts only the comparator; no
unconstrained minimizer is assumed to exist.

The expectation convention does not enter the quantitative proof. All
upper-bound protocols used below have bounded queries and bounded
sample-path regret. The lower-bound adversaries produce regret with a
uniformly bounded negative part for every protocol, including those
with nonintegrable action norms.

\subsection{The minimax characterization}

\begin{theorem}[Optimal communication--regret tradeoff]
\label{thm:main}
In the model \eqref{eq:model}, for every $M,K,R,d\ge1$ and $G,B>0$,
\begin{equation}
 c\,GB\sqrt{\frac DT}
 \le \calR_{M,K,R,d}(G,B)
 \le C\,GB\sqrt{\frac DT},
 \qquad D=\min\{T,d,K+d/M\},
 \label{eq:main}
\end{equation}
where $c,C>0$ are universal. One may take $c=1/256$ and $C=6$.
Equivalently,
\begin{equation}
 \calR_{M,K,R,d}(G,B)
 =\Theta\!\left(GB\min\left\{1,\sqrt{\frac dT},
                     \sqrt{\frac1R+\frac d{MT}}\right\}\right).
 \label{eq:main-envelope}
\end{equation}
If $G=0$ or $B=0$, the minimax value is zero.
\end{theorem}

The upper bound selects among three simple protocols using the known
parameters. The lower bound holds for every randomized protocol in the
model, with unrestricted exact-real queries and no structural
restriction on its iterates. No logarithmic factors or dimension
thresholds are hidden in \eqref{eq:main-envelope}.

\begin{table}[t]
\centering
\small
\renewcommand{\arraystretch}{1.4}
\begin{tabularx}{\textwidth}{@{}p{0.22\textwidth}p{0.35\textwidth}Y@{}}
\toprule
Parameter regime & Minimax average regret & Interpretation\\
\midrule
$d\le K$ & $\Theta(GB\sqrt{d/T})$
& Independent learning is optimal.\\
$K<d<MK$ & $\Theta(GB/\sqrt R)$
& Communication controls the rate.\\
$d\ge MK$ & $\Theta\!\left(GB\min\{1,\sqrt{d/(MT)}\}\right)$
& Pooled information controls the rate, with a constant cap.\\
\bottomrule
\end{tabularx}
\caption{Interpretation of Theorem~\ref{thm:main} up to universal
constants. The middle regime is empty when $M=1$. At a shared boundary
the expressions agree up to constants; the complete minimum in
\eqref{eq:main} gives the exact rate expression used in the theorem.}
\label{tab:regimes}
\end{table}

\subsection{Regimes and quantitative comparisons}

Table~\ref{tab:regimes} follows directly from \eqref{eq:main}. If
$d\le K$, then $D=d$, since $K\le T$. If $K<d<MK$, then
$K\le D\le2K$, giving $GB/\sqrt R$. If $d\ge MK$, the term
$K+d/M$ lies between $d/M$ and $2d/M$, and the rate is
$GB\min\{1,\sqrt{d/(MT)}\}$ up to constants. These are comparisons
of rates, not claims about exact switching points of particular
algorithms.

Several boundary cases help interpret the formula. With communication
after every step ($K=1$), it becomes
\begin{equation}
 \Theta\!\left(GB\min\left\{1,\frac1{\sqrt T}
                                  +\sqrt{\frac d{MT}}\right\}\right).
 \label{eq:sync}
\end{equation}
With a single worker ($M=1$), or with no useful communication before
the horizon ends ($R=1$), it becomes
$\Theta(GB\min\{1,\sqrt{d/T}\})$. In one dimension it is always
$\Theta(GB/\sqrt T)$. When $d\ge MT$, the regret is
$\Theta(GB)$. In this regime, the total observation budget is too small
for the worst-case regret to vanish.

\begin{corollary}[An improvement in the upper-bound guarantee]
\label{cor:improvement}
Let $s\ge2$ be an integer, and set $M=K=s$, $d=R=s^2$, and
$T=s^3$. Then
\[
 \calR_{s,s,s^2,s^2}(G,B)=\Theta(GBs^{-1}),
 \qquad U_{\rm prior}=\Theta(GBs^{-1/2}),
\]
where $U_{\rm prior}$ is the explicit expression in
\eqref{eq:prior-envelope}.
\end{corollary}
\begin{proof}
Here $D=\min\{s^3,s^2,2s\}=2s$. In the prior expression, the
independent term is $s^{-1/2}$, while the collaborative expression is
$s^{-3/2}+s^{-1}+s^{-1/2}$. The claims follow.
\end{proof}

\begin{corollary}[A necessary communication cost]
\label{cor:witness}
Let $M=d=K=s$, $R=s^2$, and $T=s^3$, with integer $s\ge2$. Then
\[
 \calR_{s,s,s^2,s}(G,B)=\Theta(GBs^{-1}).
\]
For the same $M,d,T$ with communication after every step, the minimax
regret is $\Theta(GBs^{-3/2})$.
\end{corollary}
\begin{proof}
With $K=s$, $D=\min\{s^3,s,s+1\}=s$. With $K=1$ and
$R=s^3$, $D=\min\{s^3,s,2\}=2$. Apply Theorem~\ref{thm:main}.
\end{proof}

Independent learning already gives the upper bound in
Corollary~\ref{cor:witness}. Its content is the matching lower bound:
the $\sqrt s$ cost of intermittent communication cannot be removed by
a different protocol in this adaptive model. This is distinct from
the improvement in an upper-bound guarantee established in
Corollary~\ref{cor:improvement}.

\section{Upper bound: batching and independent learning}
\label{sec:upper}

The optimal upper bound is an envelope of three strategies. We first analyze the collaborative strategy, which keeps one center throughout a communication block, and then give the independent and zero-action alternatives.

\subsection{Two-point smoothing and the second moment}

For $\delta>0$ and a convex $G$-Lipschitz loss $f$, define
\[
 F_\delta(x)=\E_{V\sim\Unif(\B(1))}f(x+\delta V).
\]
The smoothed function averages over the ball, while the estimator samples a direction from the sphere. This standard two-point construction~\cite{ADX10,Shamir17} has the following properties.

\begin{lemma}[Estimator identities]\label{lem:estimator}
Let $U\sim\Unif(\Sph)$ and
\[
 g(x,U)=\frac d{2\delta}
       \bigl(f(x+\delta U)-f(x-\delta U)\bigr)U.
\]
Then
\begin{align}
&f(x)\le F_\delta(x)\le f(x)+G\delta,\qquad
 \norm{\nabla F_\delta(x)}\le G, \label{eq:smoothing}\\
&\E g(x,U)=\nabla F_\delta(x),\qquad
 \E\norm{g(x,U)}^2\le 2dG^2. \label{eq:moment}
\end{align}
The statements remain valid conditionally when $f$ and $x$ are chosen
using past information, provided $U$ is fresh and independent of them.
\end{lemma}
\begin{proof}
Convexity, Jensen's inequality, and Lipschitzness give
\eqref{eq:smoothing}. Differentiation of the ball integral and the
divergence theorem give
\[
 \nabla F_\delta(x)=\frac d\delta
       \E_{U\sim\Unif(\Sph)} f(x+\delta U)U.
\]
Symmetry under $U\mapsto-U$ gives the unbiasedness identity.

For completeness, the second-moment calculation uses the following
sphere Poincar\'e inequality, with normalized surface measure:
\begin{equation}
 \E h(U)^2\le\frac1{d-1}\E\norm{\nabla_{\Sph}h(U)}^2
 \quad(d\ge2,\ \E h(U)=0).                    \label{eq:poincare}
\end{equation}
One way to see its constant is to expand a smooth $h$ in spherical
harmonics. A degree-$\ell$ harmonic is an eigenfunction of
$-\Delta_{\Sph}$ with eigenvalue $\ell(\ell+d-2)$, as follows by
writing the Euclidean Laplacian in polar coordinates. Integration by
parts gives its contribution
$\ell(\ell+d-2)\norm{h_\ell}_{L^2}^2$ to the Dirichlet integral.
The constant component is absent and every remaining eigenvalue is at
least $d-1$. Summing the orthogonal expansion proves
\eqref{eq:poincare}. Polynomial restrictions are dense, and decomposing
homogeneous polynomials into harmonic components yields the expansion;
smooth approximation extends the inequality to Lipschitz functions.

Apply \eqref{eq:poincare} to the odd, mean-zero function
\[
 h(u)=\frac{f(x+\delta u)-f(x-\delta u)}{2\delta}.
\]
It is $G$-Lipschitz as a function of $u$, so its tangential gradient has
norm at most $G$. Therefore
\[
 \E\norm{g(x,U)}^2=d^2\E h(U)^2
 \le\frac{d^2}{d-1}G^2\le2dG^2.
\]
For $d=1$, the estimator is a deterministic secant slope with absolute
value at most $G$. Conditioning on the chosen function and center
proves the conditional version.
\end{proof}

\subsection{Communication-block batching}

Let $I_r=\{(r-1)K+1,\ldots,rK\}$. All workers use the same center
$z_r\in\B(B)$ throughout $I_r$, starting from $z_1=0$. At each time and
worker, independently sample a fresh sphere direction $U_t^m$ and
simultaneously query
\[
 x_t^{m,1}=z_r+\delta U_t^m,\qquad
 x_t^{m,2}=z_r-\delta U_t^m,
 \qquad \delta=B/T.
\]
Compute the estimator $g_t^m$ of Lemma~\ref{lem:estimator}. At the next
permitted communication, form
\[
 \bar g_r=\frac1{MK}\sum_{t\in I_r,m\in[M]}g_t^m,
 \qquad
 z_{r+1}=\Pi_{\B(B)}(z_r-\eta\bar g_r).
\]
No update or message is needed inside the communication block.

Write
\[
 h_t^m=\nabla F_{t,\delta}^m(z_r),\quad
 \xi_t^m=g_t^m-h_t^m,\quad
 \bar h_r=\frac1{MK}\sum_{t\in I_r,m}h_t^m,\quad
 \bar\xi_r=\bar g_r-\bar h_r.
\]
Condition on all past information and the adversary's current functions
before revealing the fresh directions. Then $\E[\xi_t^m\mid\cdot]=0$,
and the estimator errors are conditionally independent across workers.
Errors from different times are orthogonal in $L^2$ by conditional
centering. Thus
\begin{align}
 \E\norm{\bar\xi_r}^2&\le\frac{2dG^2}{MK},&
 \E[\bar\xi_r\mid\text{block-start history}]&=0, \label{eq:block-noise}\\
 \E\norm{\bar g_r}^2
 &\le2G^2+\frac{4dG^2}{MK}.                 \label{eq:block-moment}
\end{align}
For the last inequality use $\norm{\bar h_r}\le G$ and
$\norm{a+b}^2\le2\norm a^2+2\norm b^2$.
The use of this inequality accounts for within-block adaptivity: future losses can depend on earlier estimator errors, so $\bar h_r$ and $\bar\xi_r$ need not be uncorrelated.

The projected-gradient inequality~\cite{Zinkevich03} holds pathwise and simultaneously for every
$u\in\B(B)$,
\begin{equation}
 \sum_{r=1}^R\ip{\bar g_r}{z_r-u}
 \le \frac{B^2}{2\eta}+\frac\eta2
          \sum_{r=1}^R\norm{\bar g_r}^2.    \label{eq:ogd}
\end{equation}
Indeed, expand
$\norm{z_r-\eta\bar g_r-u}^2$ and use nonexpansiveness of projection;
the squared distances telescope and $\norm{z_1-u}^2\le B^2$.

\paragraph{Charging both queries.}
The average queried loss satisfies the pointwise bound
\begin{align}
 \frac{f_t^m(z_r+\delta U_t^m)+f_t^m(z_r-\delta U_t^m)}2-f_t^m(u)
 &\le F_{t,\delta}^m(z_r)-F_{t,\delta}^m(u)+2G\delta \notag\\
 &\le \ip{h_t^m}{z_r-u}+2G\delta.         \label{eq:charged}
\end{align}
\paragraph{The random comparator.}
Let $u^\star$ be a minimizer of the realized cumulative original losses
over $\B(B)$. We may substitute this \emph{random} $u^\star$ in
\eqref{eq:ogd}, because that inequality holds for every $u$ on each
sample path. The correction from estimated to true gradients is
\begin{equation}
 \sum_r\ip{\bar h_r}{z_r-u^\star}
 =\sum_r\ip{\bar g_r}{z_r-u^\star}
 -\sum_r\ip{\bar\xi_r}{z_r}
 +\ip{\sum_r\bar\xi_r}{u^\star}.          \label{eq:comparator-correction}
\end{equation}
The middle sum has expectation zero because each $z_r$ is fixed at the start of its block. The final term need not be centered: $u^\star$ can depend on all estimator errors. Its magnitude is controlled by
\begin{equation}
 \E\!\left[\sup_{u\in\B(B)}\ip{\sum_r\bar\xi_r}{u}\right]
 = B\E\norm{\sum_r\bar\xi_r}
 \le BG\sqrt{\frac{2dR}{MK}}.             \label{eq:terminal}
\end{equation}
The bound is uniform over every minimizer on each sample path, so no measurable choice of $u^\star$ is needed: one can equivalently take the supremum over $u\in\B(B)$ before taking expectations. This terminal martingale correction preserves the minimum inside the expectation in \eqref{eq:path-regret}.

Combining \eqref{eq:block-moment}--\eqref{eq:terminal}, and taking
\[
 \eta=\frac{B}{G\sqrt{R(2+4d/(MK))}},
\]
gives the following bound for the average regret:
\begin{align}
 \operatorname{Reg}_{\rm batch}
 &\le
 GB\sqrt{\frac2R+\frac{4d}{MT}}
 +GB\sqrt{\frac{2d}{MT}}+\frac{2GB}{T} \notag\\
 &\le 6GB\sqrt{\frac1R+\frac d{MT}}.      \label{eq:batch-bound}
\end{align}

\subsection{The independent and trivial alternatives}

Each worker can instead run projected online gradient descent for all
$T$ steps with its own two-point estimates, $\delta=B/T$, and
$\eta=B/(G\sqrt{2dT})$. Sum the $M$ pathwise projection inequalities.
The initial potential divided by $MT$ is $B^2/(2\eta T)$ and the
quadratic term in expectation is at most $\eta dG^2$.
The predictable error pairings again have expectation zero.
The single common random comparator contributes at most
\[
 \frac{B}{MT}\E\norm{\sum_{t,m}\xi_t^m}
 \le GB\sqrt{\frac{2d}{MT}}.
\]
Consequently,
\begin{equation}
 \operatorname{Reg}_{\rm independent}
 \le GB\sqrt{\frac{2d}{T}}+
       GB\sqrt{\frac{2d}{MT}}+\frac{2GB}{T}
 \le5GB\sqrt{\frac dT}.                  \label{eq:ind-bound}
\end{equation}
Finally, querying the zero point twice gives regret at most $GB$
pointwise, by Lipschitzness and $\norm{u^\star}\le B$.
Selecting in advance among these three methods according to the known
parameters proves the upper bound in Theorem~\ref{thm:main} with $C=6$.
All three protocols respect the simultaneous-query and communication rules. Their queries have norm at most $B+B/T$, so their sample-path regret has a uniform finite bound under every admissible adversary.

\section{Tools for the lower bounds}\label{sec:tools}

The constructions in the next two sections share three ingredients. A radial penalty controls losses outside the comparator ball. A posterior calculation tracks exactly what adaptive real-valued measurements reveal about a spherical direction. Finally, a martingale estimate converts the energy retained in each episode into a lower bound on the expected norm of the cumulative coefficient.

\subsection{A differentiable penalty for unrestricted actions}

Set $\gamma=G/2$ and, for $r\ge0$, define
\begin{equation}
 \rho_B(r)=
 \begin{cases}
  0, &0\le r\le B,\\
  (r-B)^2/(2B), &B<r<2B,\\
  r-3B/2, &r\ge2B,
 \end{cases}
 \qquad
 \psi(x)=\gamma\rho_B(\norm{x}).          \label{eq:penalty}
\end{equation}
This function is convex and continuously differentiable, is zero on
$\B(B)$, and is $\gamma$-Lipschitz. Its gradient is zero near the
origin, so differentiability of the Euclidean norm at the origin is
irrelevant. Moreover, $\rho_B(r)\ge r-3B/2$.

For any $a\in\R^d$ with $\norm a\le\gamma$, the function
\begin{equation}
 f_a(x)=\ip a x+\psi(x)                  \label{eq:hard-loss}
\end{equation}
is differentiable, convex, and globally $G$-Lipschitz, and
\begin{equation}
 f_a(x)\ge-3GB/4\quad\text{for every }x\in\R^d.
 \label{eq:bounded-below}
\end{equation}
On the comparator ball it is exactly linear. The penalty is fixed,
public, and independent of every hidden direction. Giving its formula
to the workers only strengthens the information available to them.

\subsection{What adaptive scalar measurements reveal}

\begin{lemma}[Adaptive linear measurements]\label{lem:posterior}
Let $V\sim\Unif(\Sph)$ be independent of all initial side information, including the observer's random tape.
An observer may choose each vector $q_i\in\R^d$ as an arbitrary measurable
function of its preceding scalar measurements and side information,
then observe $\ip V{q_i}$. Let
$S=\spanop\{q_1,\ldots,q_q\}$ and let $P_S$ be orthogonal projection.
Conditionally on the resulting transcript and side information, $P_SV$ is known and
$(I-P_S)V$ is uniform on the centered sphere of radius
$\sqrt{1-\norm{P_SV}^2}$ in $S^\perp$; a zero-radius sphere means a point mass at zero. In particular,
\begin{equation}
 \E[(I-P_S)V\mid\text{transcript}]=0.     \label{eq:post-center}
\end{equation}
For $q\le d$,
\begin{equation}
 \E\norm{(I-P_S)V}^2\ge1-\frac qd.       \label{eq:post-energy}
\end{equation}
The same statements hold for batches of vectors selected before any
of their measurements are returned.
\end{lemma}
\begin{proof}
For a fixed current transcript, decompose a new vector into its component
in $S$ and its component in $S^\perp$. The former gives an already known
scalar. If the latter is nonzero, its direction is fixed before the
new residual measurement, so conditioning a uniform sphere on that
one coordinate leaves a uniform sphere in its orthogonal complement.
If it is zero, the transcript contains no new information. Induction
gives the claimed conditional law. Adaptivity adds no extra constraints:
the selected vectors are themselves functions of the previous transcript.

To prove the energy bound, first consider exactly $q$ linearly
independent measurement directions, orthogonalizing each against its
predecessors. At rank $r$, if the residual squared radius is $Q_r$, the
next removed squared coordinate has conditional expectation
$Q_r/(d-r)$. Hence
\[
 \E[Q_{r+1}\mid\text{rank-}r\text{ transcript}]
 =Q_r\frac{d-r-1}{d-r},\qquad Q_0=1.
\]
Iterating gives $\E Q_q=(d-q)/d$.
If there are redundant queries or an adaptively chosen smaller rank,
append measurable orthogonal completion queries until rank $q$ is
reached. This can only decrease the residual energy and proves
\eqref{eq:post-energy}. For a simultaneous batch, reveal the prechosen
vectors one at a time; they are also legal choices in this more
informative sequential measurement experiment.
\end{proof}

\begin{lemma}[Equivalence of residual responses and linear measurements]
\label{lem:transcript}
Suppose that, before a time step, $S$ is the span of all the observer's
previous query points in the current episode. The two responses are
values of
\[
 x\longmapsto \gamma\ip{(I-P_S)V}{x}+\psi(x)
\]
at a pair selected simultaneously. Starting with $S=\{0\}$, its whole
local value transcript is equivalent to the corresponding adaptive
linear measurements of $V$ at all its query points. In particular,
before a time step the conditional mean of $(I-P_S)V$ is zero.
\end{lemma}
\begin{proof}
Inductively, the preceding responses determine $P_SV$. For a newly
queried point $x$, subtract the known $\psi(x)$, divide by $\gamma$,
and add the known scalar $\ip{P_SV}{x}$. The result is precisely
$\ip Vx$. Conversely, knowledge of $P_SV$ and $\ip Vx$ reconstructs the
response. Apply Lemma~\ref{lem:posterior}.
Both points must be chosen before either new scalar is observed to
conclude that both charged linear losses have conditional mean zero
against the same current residual.
\end{proof}

\subsection{From episode energy to expected cumulative norm}

\begin{lemma}[A first-moment bound for vector martingales]\label{lem:mart}
Let $(X_b)_{b=1}^N$ be vector martingale differences for a filtration
$(\mathcal F_b)$, and let deterministic numbers $a_b>0$ satisfy
\[
 \norm{X_b}\le a_b,\qquad
 \E[\norm{X_b}^2\mid\mathcal F_{b-1}]\ge a_b^2/4.
\]
Then
\begin{equation}
 \E\norm{\sum_{b=1}^N X_b}
 \ge\frac1{16}\sqrt{\sum_{b=1}^Na_b^2}.  \label{eq:mart-lower}
\end{equation}
\end{lemma}
\begin{proof}
Put $S_n=\sum_{b\le n}X_b$ and $A=\sum_ba_b^2$. Orthogonality gives
\[
 \E\norm{S_N}^2\ge A/4,\qquad
 \E\norm{S_{b-1}}^2\le\sum_{i<b}a_i^2.
\]
For fixed $s=S_{b-1}$, expand $\norm{s+X_b}^4$. The term
$4\norm s^2\ip s{X_b}$ has conditional expectation zero. The other
mixed terms are at most
\[
 6a_b^2\norm s^2+4a_b^3\norm s
 \le8a_b^2\norm s^2+2a_b^4.
\]
Including $\norm{X_b}^4\le a_b^4$ yields
\[
 \E[\norm{S_b}^4\mid\mathcal F_{b-1}]
 \le\norm{S_{b-1}}^4+8a_b^2\norm{S_{b-1}}^2+3a_b^4.
\]
Summing and using the preceding second-moment upper bound gives
\[
 \E\norm{S_N}^4
 \le8\sum_ba_b^2\sum_{i<b}a_i^2+3\sum_ba_b^4
 \le4A^2.
\]
Finally, H\"older's inequality in the form
\[
 \E Z\ge\frac{(\E Z^2)^{3/2}}{(\E Z^4)^{1/2}}
 \quad (Z\ge0)
\]
with $Z=\norm{S_N}$ gives \eqref{eq:mart-lower}.
\end{proof}

\section{Communication lower bound from local information}
\label{sec:local}

During a communication block, each worker can use only its local measurements. The following adversary exploits this restriction while correlating the remaining directions across workers.

\begin{proposition}[Communication obstruction]\label{prop:local}
Every randomized protocol in \eqref{eq:model} has an admissible adaptive
adversary for which
\begin{equation}
 \operatorname{Reg}\ge
 \frac{GB}{64\sqrt2}\sqrt{\frac{\min\{K,d\}}{T}}.
 \label{eq:local-lower}
\end{equation}
\end{proposition}
\begin{proof}
Define
\[
 L_{\rm loc}=\min\{K,\lceil d/4\rceil\}.
\]
Partition each of the $R$ communication blocks into consecutive episodes
of length $L_{\rm loc}$, except for a possibly shorter final episode.
Thus no communication occurs in the interior of an episode. Write
their lengths as $\ell_1,\ldots,\ell_N$.

At the start of episode $b$, sample a fresh
$V_b\sim\Unif(\Sph)$, independent of all previous randomness and of the
algorithm's private or shared random tapes. At local time
$j\in\{1,\ldots,\ell_b\}$ let
\[
 S_{b,j}^m=\spanop\{x_{b,i}^{m,1},x_{b,i}^{m,2}:1\le i<j\},
\qquad
 a_{b,j}^m=\gamma(I-P_{S_{b,j}^m})V_b,
\]
and use
\begin{equation}
 f_{b,j}^m(x)=\ip{a_{b,j}^m}{x}+\psi(x).
 \label{eq:local-adversary}
\end{equation}
The adversary uses only its fresh direction and \emph{strictly previous}
played points. It never inspects a current query pair or current random
coin. The function belongs to the required class by
\eqref{eq:hard-loss}--\eqref{eq:bounded-below}.

Represent all algorithmic randomization, including shared randomness,
by a tape $W$ independent of the adversary's directions. For analysis
only, condition on $W$ and on every earlier episode's hidden direction.
This fixes all initial worker states for the current episode without
changing the uniform prior of $V_b$.
Because no communication occurs during the episode and
\eqref{eq:local-adversary} at worker $m$ uses only that worker's earlier
queries, its transcript is the experiment in
Lemma~\ref{lem:transcript}.

Let $\mathcal H_{b,j-1}^m$ contain this initial side information and
worker $m$'s local transcript before the current pair.
The current residual has conditional mean zero given
$\mathcal H_{b,j-1}^m$. Since the full algorithmic tape is included for
analysis, both simultaneously selected points are measurable with
respect to this sigma field. Therefore
\begin{equation}
 \E\bigl[f_{b,j}^m(x_{b,j}^{m,k})\mid\mathcal H_{b,j-1}^m\bigr]
 =\psi(x_{b,j}^{m,k})\ge0,\qquad k=1,2.
 \label{eq:nonnegative-charged}
\end{equation}
The adversarial rule itself does not use the tape.

This conclusion does not require integrability of the query norms.
For every fixed finite point, integration over the residual gives the
finite value $\psi(x)$. The lower bound
$f_{b,j}^m(x)\ge-3GB/4$ then justifies iterated expectation in the
extended sense. Thus arbitrary heavy-tailed or unbounded actions
cannot invalidate the argument.

Before local time $j$, the worker has at most $2(j-1)$ scalar
measurements. Our episode length ensures
\[
 2(j-1)\le2(\lceil d/4\rceil-1)<d/2.
\]
Lemma~\ref{lem:posterior} therefore gives, conditional on the initial
side information,
\begin{equation}
 \E\norm{(I-P_{S_{b,j}^m})V_b}^2\ge\frac12.
 \label{eq:local-residual-energy}
\end{equation}
Define the episode's aggregate linear coefficient
\[
 X_b=\sum_{j=1}^{\ell_b}\sum_{m=1}^Ma_{b,j}^m,\qquad
 a_b=\gamma M\ell_b.
\]
Then $\norm{X_b}\le a_b$. Also,
\[
 \ip{X_b}{V_b}
 =\gamma\sum_{j,m}\norm{(I-P_{S_{b,j}^m})V_b}^2.
\]
By \eqref{eq:local-residual-energy},
$\E[\ip{X_b}{V_b}\mid\text{initial side information}]\ge a_b/2$.
Since $\norm{V_b}=1$, Cauchy--Schwarz and Jensen give
\begin{equation}
 \E[\norm{X_b}^2\mid\text{initial side information}]\ge a_b^2/4.
 \label{eq:episode-energy}
\end{equation}

Each residual has expectation zero after averaging its local posterior.
Consequently,
$\E[X_b\mid W,V_1,\ldots,V_{b-1}]=0$.
Thus $X_b$ is a martingale difference for
$\mathcal F_b=\sigma(W,V_1,\ldots,V_b)$.
\emph{This centering is at episode boundaries.}
We do not assert centering of a worker's coefficient given the entire
global transcript within the episode: other workers may have learned
information it does not possess.

Because $\psi$ vanishes on the comparator ball,
\[
 \min_{\norm u\le B}\sum_{b,j,m}f_{b,j}^m(u)
 =-B\norm{\sum_bX_b}.
\]
Equations \eqref{eq:nonnegative-charged} and
\eqref{eq:mart-lower} now imply
\begin{equation}
 \operatorname{Reg}\ge
 \frac{B}{MT}\E\norm{\sum_bX_b}
 \ge\frac{GB}{32T}\sqrt{\sum_b\ell_b^2}. \label{eq:episode-regret}
\end{equation}

For integers $K\ge L\ge1$, splitting $K$ into pieces of length $L$
and a remainder gives a sum of squared lengths at least $KL/2$:
there are $\lfloor K/L\rfloor$ full pieces and
$\lfloor K/L\rfloor\ge K/(2L)$. Hence
\[
 \sum_b\ell_b^2\ge\frac{TL_{\rm loc}}2
 \ge\frac{T\min\{K,d\}}8.
\]
Substitute into \eqref{eq:episode-regret} to obtain
\eqref{eq:local-lower}.
\end{proof}

\section{Dimension lower bound from pooled information}
\label{sec:global}

The second adversary removes all previously measured directions, even those held at other workers. It therefore yields a lower bound for a stronger information model in which all past feedback is pooled after every step.

\begin{proposition}[Dimension obstruction, including saturation]
\label{prop:global}
Every randomized protocol in \eqref{eq:model} has an admissible adaptive adversary for which
\begin{equation}
 \operatorname{Reg}\ge
 \frac{GB}{64\sqrt2}
 \sqrt{\frac{\min\{T,d/M\}}{T}}.
 \label{eq:global-lower}
\end{equation}
This lower bound holds even if all past scalar observations are shared
after every time step.
\end{proposition}
\begin{proof}
Let
\[
 L_{\rm glob}=\min\{T,\lceil d/(4M)\rceil\},
\]
and partition the whole horizon into consecutive episodes of this
length and a possible shorter remainder. These episodes need not align
with the communication blocks.

For each episode $b$, draw a fresh uniform sphere direction $V_b$.
Before its local time $j$, define the \emph{global} subspace
\[
 S_{b,j}^{\rm glob}
 =\spanop\{x_{b,i}^{m,k}:i<j,\ m\in[M],\ k\in\{1,2\}\}.
\]
Use the same coefficient at every worker:
\begin{equation}
 a_{b,j}^m=\gamma(I-P_{S_{b,j}^{\rm glob}})V_b,
 \qquad
 f_{b,j}^m(x)=\ip{a_{b,j}^m}{x}+\psi(x).
 \label{eq:global-adversary}
\end{equation}
This is again a nonanticipating rule based only on past query points.

Give a hypothetical central observer all previous query pairs and
responses, as well as the algorithmic tape. Before a time step, its
information is precisely the corresponding adaptive linear-measurement
transcript of $V_b$: the response-equivalence argument of
Lemma~\ref{lem:transcript} works with the global subspace.
All $2M$ current query points are selected without any current response,
so every current charged linear term has conditional mean zero even
for this stronger observer.

At local time $j$, the number of preceding scalar measurements is at most
$2M(j-1)$, and
\[
 2M(j-1)\le2M(\lceil d/(4M)\rceil-1)<d/2.
\]
The residual-energy lower bound is therefore again $1/2$.
With $X_b=\sum_{j,m}a_{b,j}^m$ and $a_b=\gamma M\ell_b$,
the same calculations as in \eqref{eq:episode-energy} prove the
conditional second-moment bound, the bounded increment, and the
episode-boundary martingale property. Nonnegative expected queried
losses and the comparator identity consequently give
\[
 \operatorname{Reg}\ge
 \frac{GB}{32T}\sqrt{\sum_b\ell_b^2}.
\]
Finally,
\[
 \sum_b\ell_b^2\ge\frac{TL_{\rm glob}}2
 \ge\frac{T\min\{T,d/M\}}8,
\]
which proves \eqref{eq:global-lower}. In particular, when $d$ is of
order $MT$ or larger, the lower bound is of constant scale.
\end{proof}

\section{Proof of the minimax lower bound}
\label{sec:combine}

Let
\[
 A=\min\{K,d\},\qquad H=\min\{T,d/M\},\qquad
 D=\min\{T,d,K+d/M\}.
\]
We claim
\begin{equation}
 D\le2\max\{A,H\}.                       \label{eq:algebra}
\end{equation}
If $d\le K$, then $D=d=A$ since $K\le T$.
If $d>K$ and $d/M\ge T$, then $H=T\ge D$.
Otherwise $A=K$, $H=d/M$, and
$D\le K+d/M=A+H\le2\max\{A,H\}$.

For each fixed protocol, both propositions supply admissible adversaries,
so its worst-case regret is at least the larger of their lower bounds.
Taking the infimum over protocols, \eqref{eq:algebra} gives
\[
 \calR_{M,K,R,d}(G,B)
 \ge\frac{GB}{64\sqrt2}
       \sqrt{\frac{\max\{A,H\}}T}
 \ge\frac{GB}{128}\sqrt{\frac DT}.
\]
This proves the claimed lower bound, with the conservative constant $c=1/256$ stated in Theorem~\ref{thm:main}.
Together with Section~\ref{sec:upper}, this proves
Theorem~\ref{thm:main}. If $B=0$, querying $0$ twice gives zero regret
against the only comparator, and zero losses give a zero lower bound.
If $G=0$, all losses are constant in their argument and regret is
identically zero.

\section{Discussion}
\label{sec:discussion}

The minimax rate separates the roles of local learning, communication,
and total information. Low-dimensional problems can be learned
independently before a communication becomes useful. At intermediate
dimensions, the obstruction is the period during which useful
observations remain local. At larger dimensions, the total number of
scalar measurements becomes decisive, and collaboration achieves the
pooled-information scale until the constant-regret cap is reached.
The optimal upper bound needs only a choice between independent
learning and block batching, together with the zero-action safeguard.

The two lower bounds also explain why a single centralized information
argument is insufficient. Such an argument detects the $d/(MT)$ term
but misses information that exists elsewhere in the network and is
temporarily unavailable to a worker. The local construction retains
this distinction throughout each episode and establishes martingale
centering only at episode boundaries. Its continuous hidden directions
and exact transcript calculation apply directly to exact real-valued
feedback and arbitrary measurable randomized protocols.

The characterization is specific to a fixed periodic communication
schedule, simultaneous charged query pairs, unrestricted heterogeneity,
and a nonanticipating adaptive adversary. Choosing the second query
after the first response changes the conditional-centering argument.
Making other workers' within-block actions public changes the local
information structure. Likewise, oblivious or stochastic losses,
restrictions on communication bits, adaptive communication schedules,
and additional smoothness or heterogeneity assumptions define different
problems. The present lower bounds do not establish matching rates for
those variants. Determining which parts of the local-information
obstruction survive under such changes is a natural direction for
further work.

\begin{thebibliography}{9}

\bibitem{ADX10}
Alekh Agarwal, Ofer Dekel, and Lin Xiao.
\newblock Optimal algorithms for online convex optimization with
multi-point bandit feedback.
\newblock In \emph{Proceedings of the 23rd Annual Conference on Learning
Theory}, pages 28--40, 2010.
\newblock \url{https://oferdekel.github.io/pdf/2010AgarwalDeXi.pdf}.

\bibitem{DGSX12}
Ofer Dekel, Ran Gilad-Bachrach, Ohad Shamir, and Lin Xiao.
\newblock Optimal distributed online prediction using mini-batches.
\newblock \emph{Journal of Machine Learning Research}, 13(6):165--202, 2012.
\newblock \url{https://jmlr.org/papers/v13/dekel12a.html}.

\bibitem{PWSS23}
Kumar Kshitij Patel, Lingxiao Wang, Aadirupa Saha, and Nathan Srebro.
\newblock Federated online and bandit convex optimization.
\newblock In \emph{Proceedings of the 40th International Conference on
Machine Learning}, volume 202 of \emph{Proceedings of Machine Learning
Research}, pages 27439--27460, 2023.
\newblock \url{https://proceedings.mlr.press/v202/patel23a.html}.

\bibitem{Shamir17}
Ohad Shamir.
\newblock An optimal algorithm for bandit and zero-order convex
optimization with two-point feedback.
\newblock \emph{Journal of Machine Learning Research}, 18(52):1--11, 2017.
\newblock \url{https://jmlr.org/papers/v18/16-632.html}.

\bibitem{WBSS21}
Blake E. Woodworth, Brian Bullins, Ohad Shamir, and Nathan Srebro.
\newblock The min-max complexity of distributed stochastic convex
optimization with intermittent communication.
\newblock In \emph{Proceedings of the 34th Conference on Learning Theory},
volume 134 of \emph{Proceedings of Machine Learning Research},
pages 4386--4437, 2021.
\newblock \url{https://proceedings.mlr.press/v134/woodworth21a.html}.

\bibitem{Zinkevich03}
Martin Zinkevich.
\newblock Online convex programming and generalized infinitesimal
gradient ascent.
\newblock In \emph{Proceedings of the 20th International Conference on
Machine Learning}, pages 928--936, 2003.
\newblock \url{https://www.cs.cmu.edu/afs/cs.cmu.edu/Web/People/maz/publications/ICML03.pdf}.

\end{thebibliography}
\end{document}
